\chapter{Analysis of Boolean functions}
\label{chap:analysis}

This chapter covers the Fourier analysis of Boolean functions, following
O'Donnell's \emph{Analysis of Boolean Functions}. These tools are aimed at
small-depth circuit lower bounds and at property testing. So far only the
property-testing side is used elsewhere in the library: the BLR soundness
theorem is used by the Hadamard tester in the PCP development. Everything
formalized here lives in \texttt{Complexitylib.BooleanAnalysis} (the
\texttt{FourierExpansion} surface, with \texttt{FourierExpansion.Defs} and
\texttt{FourierExpansion.Internal}). The following are formalized: the parity
basis, the Fourier expansion, Plancherel and Parseval, the degree
decomposition, the BLR linearity test with soundness and local correction, the
spectral definitions and basic identities of noise stability, noise
sensitivity, the noise operator, derivatives, and influence (O'Donnell,
Chapter~1 and parts of Chapter~2), the reading of influence as the probability
of a pivotal flip, and the one-bit base case of the Bonami inequality. The
correlated-input readings of the noise quantities, monotone functions, $L^p$
norms, and hypercontractivity are not yet formalized. Hypercontractivity is the
main open direction: its tensorization step is what the level-1 inequality,
FKN, KKL, and Friedgut's junta theorem all need. The Linial--Mansour--Nisan
concentration bound for $\cc{AC^0}$ needs instead the Fourier degree bound for
decision trees and a Fourier-side restriction API.

\paragraph{Conventions.} The cube is $\mathbb{F}_2^n$, written
\texttt{Cube n = Fin n -> ZMod 2}. The encoding is $\chi(0) = 1$, $\chi(1) = -1$, so
bit $0$ is the value $+1$ in O'Donnell's $\{-1,1\}^n$ picture. ``Boolean-valued''
always means $\{-1,1\}$-valued. Noise stability, the noise operator, noise
sensitivity, the derivative, and the influences are all \emph{defined}
spectrally. Some combinatorial readings are theorems: influence is the
probability that flipping a coordinate changes a Boolean-valued function, and
the derivative is a twisted pointwise difference
(Theorems~\ref{thm:analysis-influence-derivative}
and~\ref{thm:analysis-influence-pivotal}). The correlated-input readings of the
noise quantities are planned (Proposition~\ref{prop:analysis-noise-correlated}).

\section{The cube and the Fourier basis}

\begin{definition}[Boolean functions on the cube]
  \label{def:analysis-cube}
  \lean{Complexity.BooleanAnalysis.Cube, Complexity.BooleanAnalysis.BooleanFunction, Complexity.BooleanAnalysis.BooleanFunction.innerProductSpace, Complexity.BooleanAnalysis.inner_def, Complexity.BooleanAnalysis.expect, Complexity.BooleanAnalysis.expect_unfold, Complexity.BooleanAnalysis.inner_eq_expect}
  \leanok
  $\mathrm{BooleanFunction}(n)$ is the real vector space of functions
  $f : \mathbb{F}_2^n \to \mathbb{R}$. Its inner product space structure comes
  from the uniform-measure inner product
  $\langle f, g \rangle = 2^{-n} \sum_x f(x) g(x)$. The uniform expectation
  $\mathbb{E}[f]$ is the average of $f$ over the cube, so
  $\mathbb{E}[f] = 2^{-n} \sum_x f(x)$ and
  $\langle f, g \rangle = \mathbb{E}_x[f(x) g(x)]$. The type is a
  \texttt{def} rather than an \texttt{abbrev}, so the $L^2$ norm does not clash
  with Mathlib's sup norm on function types.
\end{definition}

\begin{definition}[Parity functions]
  \label{def:analysis-parity}
  \lean{Complexity.BooleanAnalysis.chi, Complexity.BooleanAnalysis.parityFun, Complexity.BooleanAnalysis.IsBooleanValued}
  \leanok
  \uses{def:analysis-cube}
  $\chi : \mathbb{F}_2 \to \mathbb{R}$ sends $b \mapsto (-1)^b$. For
  $S \subseteq [n]$, the parity function is
  $\chi_S(x) = \prod_{i \in S} \chi(x_i)$, so $\chi_\emptyset = 1$. A function
  is \emph{Boolean-valued} if every value is $1$ or $-1$.
\end{definition}

\begin{theorem}[Orthonormality of parities]
  \label{thm:analysis-parity-orthonormal}
  \lean{Complexity.BooleanAnalysis.parityFun_orthonormal, Complexity.BooleanAnalysis.expect_parityFun, Complexity.BooleanAnalysis.parityFun_mul, Complexity.BooleanAnalysis.parityFun_add, Complexity.BooleanAnalysis.isBooleanValued_parityFun, Complexity.BooleanAnalysis.parityFun_injective}
  \leanok
  \uses{def:analysis-parity}
  $\langle \chi_S, \chi_T \rangle = [S = T]$ and $\mathbb{E}[\chi_S] = [S = \emptyset]$.
  Moreover, $\chi_S \chi_T = \chi_{S \triangle T}$ pointwise,
  $\chi_S(x + y) = \chi_S(x) \chi_S(y)$, every $\chi_S$ is Boolean-valued, and
  $S \mapsto \chi_S$ is injective.
\end{theorem}
\begin{proof}
  \leanok
  If $j \in S$, translating by the unit vector $e_j$ permutes the cube and
  negates $\chi_S$, so $\sum_x \chi_S(x) = -\sum_x \chi_S(x) = 0$.
  Orthonormality then follows from $\chi_S \chi_T = \chi_{S \triangle T}$,
  since $S \triangle T = \emptyset$ iff $S = T$.
\end{proof}

\begin{definition}[Fourier coefficients and weights]
  \label{def:analysis-fourier-coeff}
  \lean{Complexity.BooleanAnalysis.fourierCoeff, Complexity.BooleanAnalysis.fourierWeight, Complexity.BooleanAnalysis.fourierWeightAtDegree}
  \leanok
  \uses{def:analysis-parity}
  The Fourier coefficient of $f$ on $S$ is
  $\widehat{f}(S) = \langle f, \chi_S \rangle$, and its Fourier weight is
  $\widehat{f}(S)^2$. The weight at degree $k$ is
  $W^k[f] = \sum_{|S| = k} \widehat{f}(S)^2$.
\end{definition}

\begin{theorem}[Fourier expansion]
  \label{thm:analysis-fourier-expansion}
  \lean{Complexity.BooleanAnalysis.fourier_expansion, Complexity.BooleanAnalysis.parityFun_span, Complexity.BooleanAnalysis.fourier_uniqueness, Complexity.BooleanAnalysis.fourierCoeff_ext, Complexity.BooleanAnalysis.eq_iff_fourierCoeff, Complexity.BooleanAnalysis.eq_zero_iff_fourierCoeff}
  \leanok
  \uses{def:analysis-fourier-coeff}
  Every $f : \mathbb{F}_2^n \to \mathbb{R}$ satisfies
  $f(x) = \sum_S \widehat{f}(S) \chi_S(x)$, so the parities span. The
  coefficients are unique: if $f = \sum_S c_S \chi_S$ pointwise, then
  $c_S = \widehat{f}(S)$ for all $S$. In particular, two functions are equal iff
  all their Fourier coefficients agree, and $f = 0$ iff every coefficient
  vanishes.
\end{theorem}
\begin{proof}
  \leanok
  \uses{thm:analysis-parity-orthonormal}
  Substituting $\widehat{f}(S) = 2^{-n} \sum_y f(y) \chi_S(y)$ and
  $\chi_S(y)\chi_S(x) = \chi_S(x + y)$ gives
  $\sum_S \widehat{f}(S)\chi_S(x) = 2^{-n} \sum_y f(y) \sum_S \chi_S(x + y)$,
  and $\sum_S \chi_S(z)$ is $2^n$ at $z = 0$ and $0$ elsewhere. Uniqueness comes
  from taking inner products with each $\chi_T$. Two functions with the same
  coefficients have the same expansion.
\end{proof}

\begin{theorem}[Plancherel and Parseval]
  \label{thm:analysis-plancherel}
  \lean{Complexity.BooleanAnalysis.plancherel, Complexity.BooleanAnalysis.parseval, Complexity.BooleanAnalysis.norm_sq_eq_inner, Complexity.BooleanAnalysis.norm_sq_eq_sum_fourierCoeff_sq, Complexity.BooleanAnalysis.parseval_boolean, Complexity.BooleanAnalysis.abs_fourierCoeff_le_one}
  \leanok
  \uses{def:analysis-fourier-coeff}
  $\langle f, g \rangle = \sum_S \widehat{f}(S) \widehat{g}(S)$ and
  $\lVert f \rVert_2^2 = \langle f, f \rangle = \sum_S \widehat{f}(S)^2$. If $f$ is
  Boolean-valued, then $\sum_S \widehat{f}(S)^2 = 1$, so
  $|\widehat{f}(S)| \le 1$ for every $S$.
\end{theorem}
\begin{proof}
  \leanok
  \uses{thm:analysis-parity-orthonormal, thm:analysis-fourier-expansion}
  The parities are orthonormal and span, so they form an orthonormal basis, and
  Plancherel is Parseval's identity for that basis. For Boolean-valued $f$,
  $\langle f, f \rangle = \mathbb{E}[f^2] = 1$.
\end{proof}

\begin{definition}[Variance, covariance, probability, distance]
  \label{def:analysis-moments}
  \lean{Complexity.BooleanAnalysis.variance, Complexity.BooleanAnalysis.covariance, Complexity.BooleanAnalysis.indicator, Complexity.BooleanAnalysis.prob, Complexity.BooleanAnalysis.hammingDist}
  \leanok
  \uses{def:analysis-cube}
  $\mathrm{Var}[f] = \mathbb{E}[f^2] - \mathbb{E}[f]^2$ and
  $\mathrm{Cov}[f,g] = \mathbb{E}[fg] - \mathbb{E}[f]\mathbb{E}[g]$. The uniform
  probability of a predicate $P$ is $\Pr[P] = \mathbb{E}[1_P]$, where $1_P$ is
  the $0$--$1$ indicator of $P$, and the relative Hamming distance is
  $\mathrm{dist}(f,g) = \Pr_x[f(x) \neq g(x)]$.
\end{definition}

\begin{proposition}[Moments in Fourier form]
  \label{prop:analysis-moments-fourier}
  \lean{Complexity.BooleanAnalysis.expect_eq_fourierCoeff_empty, Complexity.BooleanAnalysis.variance_eq_sum_fourierCoeff_sq, Complexity.BooleanAnalysis.covariance_eq_sum_fourierCoeff, Complexity.BooleanAnalysis.variance_boolean, Complexity.BooleanAnalysis.inner_eq_one_sub_two_dist}
  \leanok
  \uses{def:analysis-moments, def:analysis-fourier-coeff}
  $\mathbb{E}[f] = \widehat{f}(\emptyset)$,
  $\mathrm{Var}[f] = \sum_{S \neq \emptyset} \widehat{f}(S)^2$, and
  $\mathrm{Cov}[f,g] = \sum_{S \neq \emptyset} \widehat{f}(S)\widehat{g}(S)$.
  For Boolean-valued $f$ and $g$, $\mathrm{Var}[f] = 1 - \mathbb{E}[f]^2$ and
  $\langle f, g \rangle = 1 - 2\,\mathrm{dist}(f,g)$.
\end{proposition}
\begin{proof}
  \leanok
  \uses{thm:analysis-plancherel}
  Take $\chi_\emptyset = 1$ in the definition of $\widehat{f}(\emptyset)$, then
  apply Plancherel and split off the empty set. For Boolean-valued functions,
  $f^2 = 1$ pointwise, and $f(x) g(x)$ is $1$ where $f$ and $g$ agree and $-1$
  where they differ.
\end{proof}

\begin{proposition}[Variance bounds]
  \label{prop:analysis-variance-boolean}
  \lean{Complexity.BooleanAnalysis.variance_nonneg, Complexity.BooleanAnalysis.variance_boolean_prob, Complexity.BooleanAnalysis.variance_boolean_nonneg, Complexity.BooleanAnalysis.variance_boolean_le_one, Complexity.BooleanAnalysis.variance_dist_bounds}
  \leanok
  \uses{def:analysis-moments, def:analysis-parity}
  Every $f$ has $\mathrm{Var}[f] \ge 0$. If $f$ is Boolean-valued, then
  $\mathrm{Var}[f] = 4 \Pr[f = 1] \Pr[f = -1]$ and
  $0 \le \mathrm{Var}[f] \le 1$. Moreover, with
  $\varepsilon = \min(\mathrm{dist}(f, 1), \mathrm{dist}(f, -1))$, the distance
  to the nearer constant function,
  $2\varepsilon \le \mathrm{Var}[f] \le 4\varepsilon$.
\end{proposition}
\begin{proof}
  \leanok
  \uses{prop:analysis-moments-fourier}
  The variance is a sum of squared Fourier coefficients. For Boolean-valued $f$,
  let $p = \Pr[f = 1]$ and $q = \Pr[f = -1]$. Then $p + q = 1$ and
  $\mathbb{E}[f] = p - q$, so $\mathrm{Var}[f] = 1 - (p - q)^2 = 4pq$. Also
  $\mathrm{dist}(f, 1) = q$ and $\mathrm{dist}(f, -1) = p$, and
  $2\min(p,q) \le 4pq \le 4\min(p,q)$.
\end{proof}

\begin{definition}[Densities and convolution]
  \label{def:analysis-convolution}
  \lean{Complexity.BooleanAnalysis.convolution, Complexity.BooleanAnalysis.IsDensity, Complexity.BooleanAnalysis.setDensity}
  \leanok
  \uses{def:analysis-cube}
  The convolution is $(f * g)(x) = \mathbb{E}_y[f(y)\,g(x + y)]$. A density is a
  nonnegative $\varphi$ with $\mathbb{E}[\varphi] = 1$. The density of a nonempty
  set $A$ is $\varphi_A = 1_A / \mathbb{E}[1_A]$. For $A = \emptyset$ the Lean
  definition divides by zero and returns the zero function, so theorems about
  $\varphi_A$ assume $A$ is nonempty.
\end{definition}

\begin{theorem}[Convolution theorem]
  \label{thm:analysis-convolution}
  \lean{Complexity.BooleanAnalysis.fourierCoeff_convolution, Complexity.BooleanAnalysis.convolution_density_isDensity, Complexity.BooleanAnalysis.setDensity_isDensity, Complexity.BooleanAnalysis.fourierCoeff_setDensity, Complexity.BooleanAnalysis.setDensity_singleton_zero, Complexity.BooleanAnalysis.fourierCoeff_setDensity_singleton_zero}
  \leanok
  \uses{def:analysis-convolution, def:analysis-fourier-coeff}
  $\widehat{f * g}(S) = \widehat{f}(S)\,\widehat{g}(S)$. The convolution of two
  densities is a density. For nonempty $A$, $\varphi_A$ is a density and
  $\widehat{\varphi_A}(S) = |A|^{-1} \sum_{x \in A} \chi_S(x)$. In particular,
  $\varphi_{\{0\}} = \sum_S \chi_S$, and every Fourier coefficient of
  $\varphi_{\{0\}}$ is $1$.
\end{theorem}
\begin{proof}
  \leanok
  \uses{thm:analysis-parity-orthonormal}
  Substitute $x = y + z$ and use $\chi_S(y + z) = \chi_S(y)\chi_S(z)$. The
  density claims follow by exchanging the two averages in
  $\mathbb{E}[\varphi * \psi]$ and by computing $\mathbb{E}[1_A] = |A|/2^n$.
\end{proof}

\section{Degree structure}

\begin{definition}[Degree parts and degree]
  \label{def:analysis-degree-part}
  \lean{Complexity.BooleanAnalysis.degreePart, Complexity.BooleanAnalysis.lowDegreePart, Complexity.BooleanAnalysis.fourierWeightAbove, Complexity.BooleanAnalysis.degree, Complexity.BooleanAnalysis.spectralSample}
  \leanok
  \uses{def:analysis-fourier-coeff}
  The degree-$k$ part of $f$ is $f^{=k} = \sum_{|S| = k} \widehat{f}(S) \chi_S$,
  and $f^{\le k} = \sum_{|S| \le k} \widehat{f}(S) \chi_S$. The weight above
  degree $k$ is $W^{>k}[f] = \sum_{|S| > k} \widehat{f}(S)^2$. The degree of $f$ is
  $\max\{|S| : \widehat{f}(S) \neq 0\}$, valued in $\mathbb{N} \cup \{\bot\}$ with
  $\bot$ for $f = 0$. For Boolean-valued $f$, the spectral sample is the
  probability distribution $S \mapsto \widehat{f}(S)^2$ on subsets of $[n]$.
\end{definition}

\begin{theorem}[Degree decomposition]
  \label{thm:analysis-degree-decomposition}
  \lean{Complexity.BooleanAnalysis.sum_degreePart, Complexity.BooleanAnalysis.degreePart_inner_eq_zero, Complexity.BooleanAnalysis.norm_sq_degreePart, Complexity.BooleanAnalysis.sum_fourierWeightAtDegree, Complexity.BooleanAnalysis.sum_fourierWeightAtDegree_boolean, Complexity.BooleanAnalysis.lowDegreePart_eq_sum, Complexity.BooleanAnalysis.lowDegreePart_self_inner, Complexity.BooleanAnalysis.fourierCoeff_degreePart, Complexity.BooleanAnalysis.degreePart_zero, Complexity.BooleanAnalysis.degree_le, Complexity.BooleanAnalysis.variance_eq_sum_fourierWeightAtDegree_pos}
  \leanok
  \uses{def:analysis-degree-part, def:analysis-moments}
  $f = \sum_{k=0}^{n} f^{=k}$ and
  $\widehat{f^{=k}}(T) = [|T| = k]\,\widehat{f}(T)$. For $j \neq k$ the parts
  $f^{=j}$ and $g^{=k}$ are orthogonal. Also $\lVert f^{=k} \rVert_2^2 = W^k[f]$,
  $\sum_{k=0}^n W^k[f] = \langle f, f \rangle$ (which is $1$ for Boolean-valued
  $f$), $f^{\le k} = \sum_{j \le k} f^{=j}$, and
  $\langle f^{\le k}, f^{\le k} \rangle = \sum_{j \le k} W^j[f]$. Finally,
  $f^{=0}$ is the constant $\mathbb{E}[f]$,
  $\mathrm{Var}[f] = \sum_{k=1}^n W^k[f]$, and $\deg f \le n$.
\end{theorem}
\begin{proof}
  \leanok
  \uses{thm:analysis-fourier-expansion, thm:analysis-plancherel, prop:analysis-moments-fourier}
  Regroup the Fourier expansion and Parseval's sum by $|S|$. The projection $f^{=k}$
  keeps exactly the coefficients with $|S| = k$.
\end{proof}

\section{Linearity testing}

\begin{definition}[Linear functions and the BLR test]
  \label{def:analysis-linear}
  \lean{Complexity.BooleanAnalysis.IsLinear, Complexity.BooleanAnalysis.IsMultiplicative, Complexity.BooleanAnalysis.IsClose, Complexity.BooleanAnalysis.IsCloseToProperty, Complexity.BooleanAnalysis.blrAcceptProb}
  \leanok
  \uses{def:analysis-parity, def:analysis-moments}
  $f$ is \emph{linear} if $f = \chi_S$ for some $S$. It is \emph{multiplicative}
  if $f(x + y) = f(x) f(y)$ for all $x, y$. It is $\varepsilon$-close to $g$ if
  $\mathrm{dist}(f, g) \le \varepsilon$, and $\varepsilon$-close to a property
  $P$ if it is $\varepsilon$-close to some $g$ satisfying $P$. The BLR
  acceptance probability is $\Pr_{x,y}[f(x) f(y) = f(x + y)]$ for independent
  uniform $x$ and $y$.
\end{definition}

\begin{theorem}[Linear iff multiplicative]
  \label{thm:analysis-linear-iff-multiplicative}
  \lean{Complexity.BooleanAnalysis.isLinear_iff_isMultiplicative}
  \leanok
  \uses{def:analysis-linear}
  A Boolean-valued $f$ is linear if and only if it is multiplicative.
\end{theorem}
\begin{proof}
  \leanok
  A multiplicative $\pm 1$-valued function is a character of $\mathbb{F}_2^n$.
  Take $S = \{i : f(e_i) = -1\}$ and write each $x$ as a sum of unit vectors.
\end{proof}

\begin{proposition}[BLR acceptance in Fourier form]
  \label{prop:analysis-blr-accept}
  \lean{Complexity.BooleanAnalysis.blrAcceptProb_eq, Complexity.BooleanAnalysis.blr_completeness}
  \leanok
  \uses{def:analysis-linear, def:analysis-fourier-coeff}
  For Boolean-valued $f$, the BLR test accepts with probability
  $\frac{1}{2} + \frac{1}{2} \sum_S \widehat{f}(S)^3$. A linear $f$ is accepted with
  probability $1$.
\end{proposition}
\begin{proof}
  \leanok
  \uses{thm:analysis-convolution, thm:analysis-plancherel}
  The acceptance probability is $\frac{1}{2} + \frac{1}{2} \langle f, f * f \rangle$.
  Then apply the convolution theorem.
\end{proof}

\begin{theorem}[BLR soundness]
  \label{thm:analysis-blr-soundness}
  \lean{Complexity.BooleanAnalysis.blr_soundness}
  \leanok
  \uses{def:analysis-linear}
  If a Boolean-valued $f$ passes the BLR test with probability at least
  $1 - \varepsilon$, then $f$ is $\varepsilon$-close to linear.
\end{theorem}
\begin{proof}
  \leanok
  \uses{prop:analysis-blr-accept, thm:analysis-plancherel, prop:analysis-moments-fourier}
  Bound $\sum_S \widehat{f}(S)^3 \le \max_S \widehat{f}(S)$ using Parseval, and
  convert the largest coefficient to a distance with
  $\langle f, \chi_S \rangle = 1 - 2\,\mathrm{dist}(f, \chi_S)$.
\end{proof}

\begin{proposition}[Local correction]
  \label{prop:analysis-local-correction}
  \lean{Complexity.BooleanAnalysis.local_correctability}
  \leanok
  \uses{def:analysis-linear}
  If a real-valued $f$ on $\{0,1\}^n$ is $\varepsilon$-close to $\chi_S$ (it
  differs from $\chi_S$ on at most an $\varepsilon$ fraction of inputs), then for
  every $x$, $\Pr_y[f(y) f(x + y) = \chi_S(x)] \ge 1 - 2\varepsilon$. No
  Boolean-valuedness is needed.
\end{proposition}
\begin{proof}
  \leanok
  $y$ and $x + y$ are each uniform. Apply a union bound to the two events where
  $f$ disagrees with $\chi_S$.
\end{proof}

\section{Noise}

\begin{definition}[Noise stability, noise operator, noise sensitivity]
  \label{def:analysis-noise}
  \lean{Complexity.BooleanAnalysis.noiseStability, Complexity.BooleanAnalysis.noiseStabilityBilin, Complexity.BooleanAnalysis.noiseOp, Complexity.BooleanAnalysis.noiseSensitivity}
  \leanok
  \uses{def:analysis-fourier-coeff}
  All four are defined spectrally, for $\rho \in \mathbb{R}$:
  \begin{itemize}
    \item $\mathrm{Stab}_\rho[f] = \sum_S \rho^{|S|} \widehat{f}(S)^2$;
    \item $\mathrm{Stab}_\rho[f,g] = \sum_S \rho^{|S|} \widehat{f}(S)\widehat{g}(S)$;
    \item the noise (Bonami--Beckner) operator
      $T_\rho f = \sum_S \rho^{|S|} \widehat{f}(S) \chi_S$;
    \item $\mathrm{NS}_\rho[f] = (\langle f, f \rangle - \mathrm{Stab}_\rho[f])/2$.
  \end{itemize}
\end{definition}

\begin{theorem}[The noise operator]
  \label{thm:analysis-noise-op}
  \lean{Complexity.BooleanAnalysis.fourierCoeff_noiseOp, Complexity.BooleanAnalysis.noiseOp_one, Complexity.BooleanAnalysis.noiseOp_zero, Complexity.BooleanAnalysis.noiseStability_eq_inner, Complexity.BooleanAnalysis.noiseStabilityBilin_eq_inner, Complexity.BooleanAnalysis.noiseStabilityBilin_symm, Complexity.BooleanAnalysis.noiseStabilityBilin_self, Complexity.BooleanAnalysis.noiseStabilityBilin_one}
  \leanok
  \uses{def:analysis-noise}
  $\widehat{T_\rho f}(S) = \rho^{|S|} \widehat{f}(S)$, $T_1 f = f$, and
  $T_0 f = \mathbb{E}[f] \cdot 1$. Also
  $\mathrm{Stab}_\rho[f] = \langle f, T_\rho f \rangle$ and
  $\mathrm{Stab}_\rho[f,g] = \langle f, T_\rho g \rangle$. The bilinear form is
  symmetric, restricts to $\mathrm{Stab}_\rho[f]$ on the diagonal, and equals
  $\langle f, g \rangle$ at $\rho = 1$. Combining symmetry with
  $\mathrm{Stab}_\rho[f,g] = \langle f, T_\rho g \rangle$ shows that $T_\rho$ is
  self-adjoint; this corollary is not stated as a separate declaration.
\end{theorem}
\begin{proof}
  \leanok
  \uses{thm:analysis-parity-orthonormal, thm:analysis-plancherel, thm:analysis-fourier-expansion}
  Read off coefficients using orthonormality, then apply Plancherel. The
  identities for $T_1$ and $T_0$ compare Fourier coefficients.
\end{proof}

\begin{proposition}[Basic properties of noise stability]
  \label{prop:analysis-noise-stability}
  \lean{Complexity.BooleanAnalysis.noiseStability_eq_sum_weight, Complexity.BooleanAnalysis.noiseStability_one, Complexity.BooleanAnalysis.noiseStability_zero_eq_expect_sq, Complexity.BooleanAnalysis.noiseStability_nonneg, Complexity.BooleanAnalysis.sq_expect_le_noiseStability, Complexity.BooleanAnalysis.noiseStability_mono, Complexity.BooleanAnalysis.noiseStability_le_self_inner, Complexity.BooleanAnalysis.noiseStability_boolean_le_one, Complexity.BooleanAnalysis.noiseStability_add_two_noiseSensitivity, Complexity.BooleanAnalysis.noiseSensitivity_antitone, Complexity.BooleanAnalysis.noiseSensitivity_nonneg, Complexity.BooleanAnalysis.noiseSensitivity_one}
  \leanok
  \uses{def:analysis-noise, def:analysis-degree-part}
  $\mathrm{Stab}_\rho[f] = \sum_{k=0}^n \rho^k W^k[f]$. Moreover
  $\mathrm{Stab}_1[f] = \langle f, f \rangle$,
  $\mathrm{Stab}_0[f] = \mathbb{E}[f]^2$, and $\mathrm{NS}_1[f] = 0$. For
  $\rho \ge 0$, $0 \le \mathbb{E}[f]^2 \le \mathrm{Stab}_\rho[f]$. On
  $0 \le \rho_1 \le \rho_2$,
  $\mathrm{Stab}_{\rho_1}[f] \le \mathrm{Stab}_{\rho_2}[f]$ and
  $\mathrm{NS}_{\rho_2}[f] \le \mathrm{NS}_{\rho_1}[f]$. For
  $\rho \in [0,1]$, $\mathrm{Stab}_\rho[f] \le \langle f, f \rangle$ (so
  $\mathrm{Stab}_\rho[f] \le 1$ for Boolean-valued $f$) and
  $\mathrm{NS}_\rho[f] \ge 0$. Always,
  $\mathrm{Stab}_\rho[f] + 2\,\mathrm{NS}_\rho[f] = \langle f, f \rangle$.
\end{proposition}
\begin{proof}
  \leanok
  \uses{thm:analysis-plancherel, prop:analysis-moments-fourier}
  Term-by-term comparison of the spectral sums.
\end{proof}

\begin{proposition}[Correlated-input form of noise]
  \label{prop:analysis-noise-correlated}
  \uses{def:analysis-noise}
  \notready
  Let $y \sim N_\rho(x)$ mean that each $\chi(y_i)$ independently equals
  $\chi(x_i)$ with probability $\frac{1+\rho}{2}$ and is flipped otherwise, for
  $\rho \in [-1,1]$. Then $T_\rho f(x) = \mathbb{E}_{y \sim N_\rho(x)}[f(y)]$
  and $\mathrm{Stab}_\rho[f] = \mathbb{E}_{x,\, y \sim N_\rho(x)}[f(x) f(y)]$.
  In particular, for Boolean-valued $f$ and $\delta \in [0,1]$,
  $\mathrm{NS}_{1-2\delta}[f] = \Pr[f(x) \neq f(y)]$ when each bit is flipped
  independently with probability $\delta$.
\end{proposition}
\begin{proof}
  \uses{thm:analysis-noise-op, thm:analysis-fourier-expansion}
  Compute $\mathbb{E}_{y \sim N_\rho(x)}[\chi_S(y)] = \rho^{|S|}\chi_S(x)$ and
  extend by linearity. The first step is to define the product distribution
  $N_\rho(x)$ as a finite weighted average.
\end{proof}

\section{Influence and derivatives}

\begin{definition}[Influence and total influence]
  \label{def:analysis-influence}
  \lean{Complexity.BooleanAnalysis.influence, Complexity.BooleanAnalysis.totalInfluence}
  \leanok
  \uses{def:analysis-fourier-coeff}
  Spectrally, $\mathrm{Inf}_i[f] = \sum_{S \ni i} \widehat{f}(S)^2$ and
  $\mathbf{I}[f] = \sum_S |S|\,\widehat{f}(S)^2$.
\end{definition}

\begin{proposition}[Basic properties of influence]
  \label{prop:analysis-influence-basic}
  \lean{Complexity.BooleanAnalysis.totalInfluence_eq_sum_influence, Complexity.BooleanAnalysis.totalInfluence_eq_sum_weight, Complexity.BooleanAnalysis.influence_nonneg, Complexity.BooleanAnalysis.totalInfluence_nonneg, Complexity.BooleanAnalysis.influence_le_totalInfluence, Complexity.BooleanAnalysis.influence_le_self_inner, Complexity.BooleanAnalysis.influence_boolean_le_one, Complexity.BooleanAnalysis.totalInfluence_le, Complexity.BooleanAnalysis.totalInfluence_boolean_le, Complexity.BooleanAnalysis.totalInfluence_le_of_degree, Complexity.BooleanAnalysis.totalInfluence_boolean_le_of_degree, Complexity.BooleanAnalysis.totalInfluence_parityFun, Complexity.BooleanAnalysis.influence_parityFun}
  \leanok
  \uses{def:analysis-influence, def:analysis-degree-part}
  $\mathbf{I}[f] = \sum_i \mathrm{Inf}_i[f] = \sum_{k=0}^n k\,W^k[f]$. Each
  influence satisfies $0 \le \mathrm{Inf}_i[f] \le \mathbf{I}[f]$ and
  $\mathrm{Inf}_i[f] \le \langle f, f \rangle$, which is $1$ for Boolean-valued
  $f$. Also $\mathbf{I}[f] \le n \langle f, f \rangle$, which is $n$ for
  Boolean-valued $f$. If every coefficient on sets of size greater than $d$
  vanishes, then $\mathbf{I}[f] \le d \langle f, f\rangle$, hence
  $\mathbf{I}[f] \le d$ for Boolean-valued $f$. For parities,
  $\mathbf{I}[\chi_S] = |S|$ and $\mathrm{Inf}_i[\chi_S] = [i \in S]$.
\end{proposition}
\begin{proof}
  \leanok
  \uses{thm:analysis-plancherel, thm:analysis-parity-orthonormal}
  Exchange the sums over $i$ and $S$, and bound $|S|$ by $n$ or $d$.
\end{proof}

\begin{definition}[Coordinate flip, derivative, sensitivity operator]
  \label{def:analysis-derivative}
  \lean{Complexity.BooleanAnalysis.flipCoord, Complexity.BooleanAnalysis.flipEquiv, Complexity.BooleanAnalysis.flipCoord_flipCoord, Complexity.BooleanAnalysis.derivative, Complexity.BooleanAnalysis.sensitivityOp}
  \leanok
  \uses{def:analysis-fourier-coeff}
  $x^{\oplus i}$ flips coordinate $i$ of $x$. This is an involutive equivalence of
  the cube. The derivative is defined spectrally:
  $D_i f = \sum_{S \ni i} \widehat{f}(S)\, \chi_{S \setminus \{i\}}$. The
  sensitivity operator is $L_i f(x) = (f(x) - f(x^{\oplus i}))/2$. These are
  different operators: by Theorem~\ref{thm:analysis-influence-derivative},
  $D_i$ removes $i$ from each frequency, while $L_i$ keeps only the frequencies
  that contain $i$.
\end{definition}

\begin{theorem}[Influence as a squared norm]
  \label{thm:analysis-influence-derivative}
  \lean{Complexity.BooleanAnalysis.fourierCoeff_derivative, Complexity.BooleanAnalysis.influence_eq_norm_sq_derivative, Complexity.BooleanAnalysis.fourierCoeff_sensitivityOp, Complexity.BooleanAnalysis.influence_eq_norm_sq_sensitivityOp, Complexity.BooleanAnalysis.sensitivityOp_eq_parityFun_mul_derivative, Complexity.BooleanAnalysis.fourierCoeff_comp_flipCoord, Complexity.BooleanAnalysis.parityFun_flipCoord, Complexity.BooleanAnalysis.expect_flipCoord}
  \leanok
  \uses{def:analysis-derivative, def:analysis-influence}
  The flip preserves expectations, $\mathbb{E}_x[g(x^{\oplus i})] = \mathbb{E}[g]$,
  and $\chi_S(x^{\oplus i}) = (-1)^{[i \in S]} \chi_S(x)$. Precomposing with the
  flip gives
  $\widehat{f \circ \oplus_i}(T) = (-1)^{[i \in T]} \widehat{f}(T)$. Moreover
  $\widehat{D_i f}(T) = [i \notin T]\,\widehat{f}(T \cup \{i\})$ and
  $\widehat{L_i f}(T) = [i \in T]\,\widehat{f}(T)$. Hence
  $\mathrm{Inf}_i[f] = \lVert D_i f \rVert_2^2 = \lVert L_i f \rVert_2^2$ and
  $L_i f = \chi_{\{i\}} \cdot D_i f$.
\end{theorem}
\begin{proof}
  \leanok
  \uses{thm:analysis-plancherel, thm:analysis-parity-orthonormal, thm:analysis-fourier-expansion}
  The flip is measure-preserving and negates exactly the parities containing $i$.
  Then apply Parseval, and compare Fourier coefficients for
  $L_i f = \chi_{\{i\}} \cdot D_i f$.
\end{proof}

\begin{theorem}[Influence as pivotality]
  \label{thm:analysis-influence-pivotal}
  \lean{Complexity.BooleanAnalysis.influence_boolean_eq_expect_sensitive, Complexity.BooleanAnalysis.totalInfluence_boolean_eq_expect_sensitive}
  \leanok
  \uses{def:analysis-influence, def:analysis-derivative}
  For Boolean-valued $f$,
  $\mathrm{Inf}_i[f] = \Pr_x[f(x) \neq f(x^{\oplus i})]$ and
  $\mathbf{I}[f] = \mathbb{E}_x\bigl[\#\{i : f(x) \neq f(x^{\oplus i})\}\bigr]$.
  The Lean statements write the probability as the expectation of the
  $0$--$1$ indicator.
\end{theorem}
\begin{proof}
  \leanok
  \uses{thm:analysis-influence-derivative}
  For Boolean-valued $f$, $(L_i f)^2$ is the indicator that $i$ is pivotal at $x$.
\end{proof}

\begin{theorem}[Poincar\'e inequality]
  \label{thm:analysis-poincare}
  \lean{Complexity.BooleanAnalysis.variance_le_totalInfluence}
  \leanok
  \uses{def:analysis-influence, def:analysis-moments}
  $\mathrm{Var}[f] \le \mathbf{I}[f]$ for every $f : \mathbb{F}_2^n \to \mathbb{R}$.
\end{theorem}
\begin{proof}
  \leanok
  \uses{prop:analysis-moments-fourier}
  Every nonempty $S$ has $|S| \ge 1$.
\end{proof}

\begin{proposition}[Noise sensitivity is bounded by influence]
  \label{prop:analysis-noise-sensitivity-influence}
  \lean{Complexity.BooleanAnalysis.noiseStability_one_sub_le, Complexity.BooleanAnalysis.noiseSensitivity_le_influence}
  \leanok
  \uses{def:analysis-noise, def:analysis-influence}
  For $\rho \in [0,1]$,
  $\langle f, f \rangle - \mathrm{Stab}_\rho[f] \le (1 - \rho)\,\mathbf{I}[f]$.
  Equivalently, $\mathrm{NS}_\rho[f] \le \frac{1-\rho}{2}\,\mathbf{I}[f]$.
\end{proposition}
\begin{proof}
  \leanok
  \uses{thm:analysis-plancherel}
  Use $1 - \rho^k \le (1 - \rho) k$ term by term.
\end{proof}

\begin{proposition}[Elementary level-1 bounds]
  \label{prop:analysis-level1-elementary}
  \lean{Complexity.BooleanAnalysis.fourierCoeff_singleton_sq_le_influence, Complexity.BooleanAnalysis.fourierWeightAtDegree_one_le_totalInfluence, Complexity.BooleanAnalysis.sum_fourierCoeff_singleton_sq_le_totalInfluence, Complexity.BooleanAnalysis.sum_fourierCoeff_singleton_sq_eq_fourierWeightAtDegree_one}
  \leanok
  \uses{def:analysis-influence, def:analysis-fourier-coeff}
  $\widehat{f}(\{i\})^2 \le \mathrm{Inf}_i[f]$ and
  $\sum_i \widehat{f}(\{i\})^2 = W^1[f] \le \mathbf{I}[f]$. These are not the
  hypercontractive level-1 inequality of
  Theorem~\ref{thm:analysis-level1-inequality}.
\end{proposition}
\begin{proof}
  \leanok
  The singletons are exactly the sets of size $1$. Each inequality drops
  nonnegative terms from a spectral sum.
\end{proof}

\section{Monotone functions}

\begin{definition}[Monotone functions]
  \label{def:analysis-monotone}
  \uses{def:analysis-parity}
  Order the cube coordinatewise by the values $\chi(x_i) \in \{-1, 1\}$. So
  $x \le y$ iff $y$ is obtained from $x$ by changing some coordinates from $1$ to
  $0$. A function $f$ is \emph{monotone} if $x \le y$ implies $f(x) \le f(y)$.
  This is O'Donnell's order on $\{-1,1\}^n$. It is the reverse of the bitwise
  order on $\mathbb{F}_2^n$, and the Lean definition must record the choice
  explicitly. The circuit library's \texttt{IsMonotoneBoolFun} is a different
  notion: monotonicity of $\{0,1\}$-valued functions on bit strings in the
  bitwise order.
\end{definition}

\begin{theorem}[Influences of monotone functions]
  \label{thm:analysis-monotone-influence}
  \uses{def:analysis-monotone, def:analysis-influence}
  \notready
  If $f$ is Boolean-valued and monotone, then
  $\mathrm{Inf}_i[f] = \widehat{f}(\{i\})$ for every $i$, and
  $\mathbf{I}[f] \le \sqrt{n}$.
\end{theorem}
\begin{proof}
  \uses{thm:analysis-influence-pivotal, thm:analysis-plancherel}
  $\widehat{f}(\{i\}) = \mathbb{E}[f(x)\chi(x_i)]$. By monotonicity this equals
  the probability that $i$ is pivotal. For the second claim, apply Cauchy--Schwarz
  to $\sum_i \widehat{f}(\{i\})$ and use $W^1[f] \le 1$.
\end{proof}

\begin{theorem}[Margulis--Russo formula]
  \label{thm:analysis-margulis-russo}
  \uses{def:analysis-monotone, def:analysis-influence}
  \notready
  Let $\mu_p$ be the product measure on the cube in which each coordinate
  independently has $\chi(x_i) = 1$ with probability $p$. For monotone
  Boolean-valued $f$,
  $\frac{d}{dp} \mu_p(f = 1) = \sum_i \Pr_{x \sim \mu_p}[f(x) \neq f(x^{\oplus i})]$.
\end{theorem}
\begin{proof}
  \uses{thm:analysis-influence-pivotal}
  Differentiate the multilinear polynomial $p \mapsto \mu_p(f = 1)$ one coordinate
  at a time. The first step is a $p$-biased expectation, which the library does
  not yet have (it has only the uniform measure).
\end{proof}

\section{Hypercontractivity}

\begin{lemma}[Bonami base case]
  \label{lem:analysis-bonami-base}
  \lean{Complexity.BooleanAnalysis.two_point_hypercontractive}
  \leanok
  For all real $a, b$,
  $\frac{1}{2}\bigl((a + b/\sqrt{3})^4 + (a - b/\sqrt{3})^4\bigr) \le (a^2 + b^2)^2$.
  Read on one bit, this says
  $\lVert T_{1/\sqrt{3}} f \rVert_4^4 \le \lVert f \rVert_2^4$ for
  $f(x) = a + bx$. The Lean statement is only the real inequality; the library
  has no $L^4$ norm or one-bit noise operator.
\end{lemma}
\begin{proof}
  \leanok
  Expand. The slack is $8b^4/9 \ge 0$.
\end{proof}

\begin{definition}[$L^p$ norms]
  \label{def:analysis-lp-norm}
  \uses{def:analysis-cube}
  For real $p \ge 1$, $\lVert f \rVert_p = \mathbb{E}[|f|^p]^{1/p}$ under the
  uniform measure. It should agree with the existing $L^2$ norm at $p = 2$ and
  satisfy H\"older's inequality.
\end{definition}

\begin{theorem}[$(2,4)$-hypercontractivity]
  \label{thm:analysis-hypercontractivity}
  \uses{def:analysis-lp-norm, def:analysis-noise}
  \notready
  For every $n$ and every $f : \mathbb{F}_2^n \to \mathbb{R}$,
  $\lVert T_{1/\sqrt{3}} f \rVert_4 \le \lVert f \rVert_2$.
\end{theorem}
\begin{proof}
  \uses{lem:analysis-bonami-base, thm:analysis-noise-op}
  Induct on $n$. Write $f = g + x_n h$, where $g$ and $h$ do not depend on
  $x_n$. Apply the base case pointwise, then the induction hypothesis to $g$ and
  $h$, with Cauchy--Schwarz. The noise operator factors as a product of one-bit
  operators, so a coordinate-restriction API for $T_\rho$ is needed first.
\end{proof}

\begin{corollary}[Bonami lemma]
  \label{cor:analysis-bonami-lemma}
  \uses{def:analysis-lp-norm, def:analysis-degree-part}
  \notready
  If $f$ has degree at most $k$, then
  $\lVert f \rVert_4 \le \sqrt{3}^{\,k}\, \lVert f \rVert_2$.
\end{corollary}
\begin{proof}
  \uses{thm:analysis-hypercontractivity, thm:analysis-noise-op}
  Apply the theorem to $T_{\sqrt{3}} f$, whose coefficients are at most
  $\sqrt{3}^{\,k}$ times those of $f$.
\end{proof}

\begin{corollary}[Dual $(4/3, 2)$-hypercontractivity]
  \label{cor:analysis-hypercontractivity-dual}
  \uses{def:analysis-lp-norm, def:analysis-noise}
  \notready
  $\lVert T_{1/\sqrt{3}} f \rVert_2 \le \lVert f \rVert_{4/3}$ for every $f$.
\end{corollary}
\begin{proof}
  \uses{thm:analysis-hypercontractivity, thm:analysis-noise-op}
  $T_{1/\sqrt{3}}$ is self-adjoint, and $4/3$ and $4$ are H\"older conjugates.
\end{proof}

\section{Consequences of hypercontractivity}

\begin{theorem}[Level-1 inequality]
  \label{thm:analysis-level1-inequality}
  \uses{def:analysis-fourier-coeff, def:analysis-cube}
  If $f : \mathbb{F}_2^n \to \{0, 1\}$ has mean $\alpha \le 1/2$, then
  $W^1[f] \le O(\alpha^2 \log(1/\alpha))$.
\end{theorem}
\begin{proof}
  \uses{thm:analysis-hypercontractivity, prop:analysis-level1-elementary}
  Let $\ell = f^{=1}$. Bound $W^1[f] = \langle f, \ell \rangle$ by splitting on
  whether $|\ell|$ exceeds a threshold. The tail is controlled by
  hypercontractive concentration of the linear form $\ell$.
\end{proof}

\begin{theorem}[Friedgut--Kalai--Naor]
  \label{thm:analysis-fkn}
  \uses{def:analysis-linear, def:analysis-fourier-coeff}
  If $f$ is Boolean-valued and $W^1[f] \ge 1 - \delta$, then $f$ is
  $O(\delta)$-close to $\chi_{\{i\}}$ or to $-\chi_{\{i\}}$ for some $i$.
\end{theorem}
\begin{proof}
  \uses{cor:analysis-bonami-lemma, thm:analysis-plancherel}
  $\ell = f^{=1}$ has $\ell^2$ close to $1$ in $L^2$. The Bonami lemma applied to
  the degree-2 polynomial $\ell^2 - 1$ forces one coefficient of $\ell$ to
  dominate.
\end{proof}

\begin{theorem}[Kahn--Kalai--Linial]
  \label{thm:analysis-kkl}
  \uses{def:analysis-influence, def:analysis-moments}
  There is an absolute constant $c > 0$ such that for every $n \ge 1$, every
  Boolean-valued $f : \mathbb{F}_2^n \to \mathbb{R}$ has a coordinate $i$ with
  $\mathrm{Inf}_i[f] \ge c \cdot \mathrm{Var}[f] \cdot \frac{\log n}{n}$.
\end{theorem}
\begin{proof}
  \uses{cor:analysis-hypercontractivity-dual, thm:analysis-influence-derivative, thm:analysis-poincare}
  Apply the $(4/3, 2)$ bound to each $D_i f$. For Boolean-valued $f$,
  $\lVert D_i f \rVert_{4/3}^{4/3} = \mathrm{Inf}_i[f]$, which gives
  $\sum_i \mathrm{Stab}_{1/3}[D_i f] \le \sum_i \mathrm{Inf}_i[f]^{3/2}$. Compare
  with the low-degree part of $\mathbf{I}[f]$.
\end{proof}

\begin{theorem}[Friedgut's junta theorem]
  \label{thm:analysis-friedgut}
  \uses{def:analysis-influence, def:analysis-linear}
  \notready
  Call $g$ a $k$-junta if it depends on at most $k$ coordinates. For every
  Boolean-valued $f$ and $0 < \varepsilon \le 1$, $f$ is $\varepsilon$-close to a
  Boolean-valued $2^{O(\mathbf{I}[f]/\varepsilon)}$-junta.
\end{theorem}
\begin{proof}
  \uses{cor:analysis-hypercontractivity-dual, thm:analysis-influence-derivative}
  Let $J$ be the coordinates with influence above
  $2^{-O(\mathbf{I}[f]/\varepsilon)}$. Hypercontractivity bounds the Fourier
  weight outside $J$ and above degree $O(\mathbf{I}[f]/\varepsilon)$. Round the
  restriction of $f$ to $J$ to $\pm 1$. The junta predicate has to be defined
  first.
\end{proof}

\section{Fourier concentration of constant-depth circuits}

These nodes connect this chapter to the circuit lower-bound chapter. The
Linial--Mansour--Nisan argument consumes the switching lemma formalized there
(random restrictions collapse small-depth circuits to shallow decision trees),
together with the fact that a shallow decision tree has low Fourier degree.
Nothing in this section is formalized yet.

\begin{definition}[Real-valued view of a Boolean function on bit strings]
  \label{def:analysis-bridge}
  \uses{def:bitstring, def:analysis-parity}
  Every $F : \{0,1\}^n \to \{0,1\}$ on bit strings determines the Boolean-valued
  function $x \mapsto \chi(F(x))$ on the cube. This must be a public bridge.
  The PCP development has only a related map, \texttt{Complexity.signOf} in an
  internal module, which sends an $\mathbb{F}_2$-valued function on the cube
  to its $\pm 1$ encoding. It does not act on bit strings.
\end{definition}

\begin{theorem}[Decision trees have low degree]
  \label{thm:analysis-decision-tree-degree}
  \uses{def:analysis-bridge, def:analysis-degree-part, def:lowerbounds-decision-tree}
  \notready
  If $F$ is computed by a decision tree of depth $t$, then the real-valued view of
  $F$ has degree at most $t$.
\end{theorem}
\begin{proof}
  \uses{thm:analysis-fourier-expansion}
  $F$ is a sum over leaves of the leaf label times the indicator of the path. The
  indicator of a depth-$t$ path is a product of $t$ functions $(1 \pm \chi(x_i))/2$.
\end{proof}

\begin{theorem}[Linial--Mansour--Nisan]
  \label{thm:analysis-lmn}
  \uses{def:analysis-bridge, def:analysis-degree-part, def:circuit, def:circuits-bases, def:NC-AC}
  \notready
  Suppose $F$ is computed by an unbounded fan-in AND/OR circuit with negations of
  size $s$ and depth $d$. Then its real-valued view satisfies
  $W^{>k}[F] \le s \cdot 2^{-\Omega(k^{1/d})}$ for every $k$.
\end{theorem}
\begin{proof}
  \uses{thm:analysis-decision-tree-degree, thm:analysis-degree-decomposition, thm:lowerbounds-switching, thm:lowerbounds-staged-switching, thm:lowerbounds-ac0-normalization}
  Relate $W^{>k}[F]$ to the expected high-degree weight of a random restriction
  of $F$. The restriction formula is
  $\widehat{F_{J \mid z}}(S) = \sum_{T \subseteq \bar J} \widehat{F}(S \cup T)\chi_T(z)$.
  By the switching lemma, the restricted circuit is a shallow decision tree with
  high probability. A Fourier-side restriction API on the cube is needed first.
\end{proof}

\begin{corollary}[Parity is not in $\cc{AC^0}$, via LMN]
  \label{cor:analysis-lmn-parity}
  \uses{def:analysis-bridge, def:circuit, def:circuits-bases, def:lowerbounds-parity}
  \notready
  For each fixed $d$, every unbounded fan-in circuit of depth $d$ computing the
  parity of $n$ bits has size $2^{\Omega(n^{1/d})}$. Håstad's bound
  (Theorem~\ref{thm:lowerbounds-hastad-size}) improves the exponent to
  $n^{1/(d-1)}$.
\end{corollary}
\begin{proof}
  \uses{thm:analysis-lmn, prop:analysis-influence-basic}
  The real-valued view of parity is $\chi_{[n]}$, whose entire Fourier weight
  sits at degree $n$.
\end{proof}
